(1) 新阱 $0<x<2a$ 的归一化本征态为 $\phi_n(x)=\sqrt{1/a}\sin(n\pi x/(2a))$，对应 $E_n'=n^2\pi^2\hbar^2/[2m(2a)^2]=n^2\pi^2\hbar^2/(8ma^2)$。突然近似的关键是刚移动墙壁时状态作为平方可积函数保持原来的形状，在 $a<x<2a$ 为零。
(2) 存活振幅 $c_1=\int_0^{2a}\phi_1(x)^*\psi_i(x)dx=\sqrt2/a\int_0^a\sin(\pi x/2a)\sin(\pi x/a)dx$。用三角恒等式积分得 $c_1=4\sqrt2/(3\pi)$，因而概率 $P_1=|c_1|^2=32/(9\pi^2)\simeq0.3603$，并非一半或一。
(3) 初始波函数在 $x=a$ 处连续归零，延拓至新阱后其导数存在有限跳变，但不存在 delta 型导数项导致无界期望能量。新哈密顿量内仍只有动能，$\langle H_f\rangle=(\hbar^2/2m)\int_0^{2a}|\psi_i'(x)|^2dx=\pi^2\hbar^2/(2ma^2)=4E_1'$。壁突然外移使状态未瞬间变化，平均动能因此保持原值。
(4) 任意新能级的系数 $c_n=\int_0^a\sqrt{2}/a\sin(\pi x/a)\sin(n\pi x/2a)dx$，演化为 $\Psi(x,t)=\sum_{n\ge1}c_n\phi_n(x)e^{-iE_n't/\hbar}$。新基组在 $L^2(0,2a)$ 完备且初态归一，故 Parseval 恒等式给 $\sum_{n\ge1}|c_n|^2=1$；瞬时只改变能量本征分解，不会使总概率丢失。
